\documentclass[10pt]{article}

\usepackage[utf8]{inputenc}

\usepackage[T1]{fontenc}

\usepackage{amsmath}

\usepackage{amsfonts}

\usepackage{amssymb}

\usepackage[version=4]{mhchem}

\usepackage{stmaryrd}

\usepackage{bbold}

\usepackage{graphicx}

\usepackage[export]{adjustbox}

\graphicspath{ {./images/} }



\begin{document}

натурального числа $n$, если в каждое конечное открытое П. этого пространства можно вписать конечное открытое П., кратность к-рого в каждой точке (т. е. число элементов П., содержащих данную точку) не превосходит $n+1$. Отношение вписанности одного П. в другое является основным общим әлементарным отношением между П. Семейство множеств $\gamma$ вписано в семейство множеств $\lambda$, если каждый элемент семейства $\gamma$ содержится в нек-ром әлементе семейства $\lambda$. На языке открытых П. определяется класс паракомпактных пространств.



П о д пок р ы т е м покрытия $\gamma$ множества $X$ наз. всякое подсемейство семейства $\gamma$, само являющееся П. множества $X$. В терминах подпокрытий определяются фундаментальные понятия бикомпактности, счетной и финальной компактности. Пространство б ик о п п а т н о, если из каждого его открытого П. можно выделить конечное подшокрытие. Пространство с ч е т н о к о м п а к т н о, если из каждого его счетного открытого П. можно выделить конечное подпокрытие. Пространство ф্ и н а л ь н о к о м п а к т в о, если из каждого его открытого П. можно выделить счетное подпокрытие. С помощью открытых П. определяются абстрактные комбинаторные объекты, открывающие дорогу применению алгебраич. методов к исследованию топологич. пространств, более общих, чем полиэдры. П. С. Александров определил фундаментальное понятие н е р в а произвольного покрытия $\gamma$ как абстрактного комплекса, вершины к-рого поставлены во взаимно однозначное соответствие с элементами покрытия $\gamma$ и конечный набор этих вершин составляет абстрактный симплекс в том и только в том случае, если пересечение отвечающих этим элементам покрытий $\gamma$ не пусто. Системам открытых П. пространства, взятых вместе с отношением вписанности, отвечают системы абстрактных комплексов, связанных симплициальными отображениями, - т. н. с п е к т ры комплекс о в.



Заметную роль в топологии играют и 3 а м к н ут ы е п о к ры т и я - то есть П., все әлементы к-рых являются замкнутыми множествами. Если в тонологич. иространстве все одноточечные множества замкнуты, то примером замкнутого П. этого пространства может служить семейство всех его одноточечных подмножеств. Но в таком П. не заключено никакой информации о топологии рассматриваемого пространства, кроме той, что в нем выполнена $T_{1}$-аксиома отделимости. Поәтому требование замкнутости П. следует соединять с другими существенными ограничениями. В частности, полезно рассматривать замкнутые локально конечные П. Они супественны в теории размерности. Важным примером замкнутого П. является П. полиэдра замкнутыми симплексами какого-нибудь подразделяющего этот полиэдр комплекса.



Среди ограничений на П., связанные не с характером әлементов, а с их расположением, наиболее часто встречаются следующие. М о н о с т ь П.- число элементов в нем, ло к а л ь н а я к о н е ч но с т ь (у каждой точки пространства есть окрестность, пересекающаяся с конечным множеством элементов семейства подмножеств этого пространства), т о ч е чн а я к о н е ч н о с т в (означающая, что множество әлементов П., содержащих произвольно взятую точку, конечно), 3 в е 3 д н а я к о н е ч н о с т в (каждый әлемент П. пересекается лишь с конечным числом әлементов этого П.).



Семейство множеств в топологич. пространстве наз. к о н с е р в а т и в ны м, если замыкание объединения любого его подсемейства равно объединению замыканий әлементов этого подсемейства. Каждое локально конечное семейство множеств консервативно. Консервативные П. возникают при исследовании пара- компактных пространств, при этом важны и не тривиальны не только открытые, но и любые консервативные П.



Важную роль играет понятие 3 в е 3 ды т о ч к $x$ относительно семейств множеств (в частности, покрытий) $\gamma$. Это - объединение всех әлементов семейства $\gamma$, содержащих $x$, обозначаемое обычно $\operatorname{St}_{\gamma}(x)$. Аналогично определяется 3 в е 3 д а $\operatorname{St}_{\gamma}(A)$ м н ож е с т в а $A$ относительно семейства множеств $\gamma$. На понятии звезды основано фундаментальное отношение звездной вписанности одного П. в другое, существенно более тонкое, чем отношение вписанности. Семейство множеств $\lambda$ наз. 3 в е 3 д н о в п и с а нн ы м в семейство множеств $\gamma$, если для каждой точки найдется элемент семейства $\gamma$, содержащий звезду этой точки относительно $\lambda$. Отношение звездной вписанности открытых П. имеет важное значение в теории размерности, на нем основаны нек-рые критерии метризуемости, оно является одним из основных элементарных понятий, входящих в определение равномерной структуры и равномерного пространства. Полезно рассматривать семейства $F$ открытых П. топологич. пространства, направленные отношением звездной вписанности в следующем смысле: для любых $\gamma_{1}$ и $\gamma_{2}$ из $F$ найдется $\mu \in F$, звездно вписанное и в $\gamma_{1}$, и в $\gamma_{2}$.



Представляет ценность следующая характеристика паракомпактности на языке звездной вписанности (т е о р е м а M о р и т ы): хаусдорфово пространство паракомпактно в том и только в том случае, если в любое его открытое П. можно вписать открытое П. звездно.



Звездная вписанность в случае произвольных (или даже замкнутых) П. не столь содержательна. В частности, это видно из того, что семейство всех одноточечных подмножеств пространства звездно вписано в любое П. этого пространства.



Лит.: [1] К е л л и Д ж., Общая топология, пер. с англ., 2 изд., М., 1981; [2] А р х а н г е ль с к и й А. В., П о н о ма р е в В. И., Основы общей топологии в задачах и упражнениях,



\begin{enumerate}

  \setcounter{enumi}{1}

  \item В ко м би на то рно й ге о м т р и и имеется ряд задач и предложений, относящихся к специальным П., в основном выпуклых множеств. Пусть $K$ - выпуклое тело $n$-мерного векторного пространства $\mathbb{R}^{n}$, bd $K$ и int $K$ - соответственно граница и внутренность $K$. Наиболее известны следующие задачи c $\Pi$

\end{enumerate}



a) Ищется минимальное число $t(K)$ транслятов (параллельных переносов) int $K$, при помощи к-рых можно Іокрыть тело $K$.



б) Ищется минимальное число $b(K)$ гомотетичных $K$ тел с коэффицциентом гомотетии $k, 0<k<1$, при гомоги к-рых можно покрыть тело $K$.



в) Ищется минимальное число $d(K)$ гомотетичных $K$ множеств с коэффициентом гомотетии $k>1$ и центром гомотетии в $\mathbb{R}^{n} \backslash$ int $K$, при помопи к-рых можно покрыть тело $K$.



При ограниченности $K$ задачи а) и б) эквивалентны между собой, эквивалентны освещения задаче (извне) множества bd $K$ и связаны с Хадвигера гипотезой. Для неограниченного $K$ задачи а) и б), вообще говоря, различны, причем числа $b(K)$ и $t(K)$ могут быть бесконечными.



\begin{center}

\includegraphics[max width=\textwidth]{2023_07_08_2fc4a688c50ef89e4641g-1}

\end{center}



\begin{center}

\includegraphics[max width=\textwidth]{2023_07_08_2fc4a688c50ef89e4641g-1(2)}

\end{center}



\includegraphics[max width=\textwidth, center]{2023_07_08_2fc4a688c50ef89e4641g-1(1)}

ние фигур на меньтие части, М., 1971; [4] Х а д в и г е р $\Gamma$, Д е б р у н не $\mathrm{p}$ Г.: Комбинаторная геометрия плоскости, пер. с нем., М., 1965: [5] Р о д ж е р с к., Укладки и покрытия, пер. с англ., М., 1968; [6] Б о Л Т я н с к и й В. Г., С о лт а н П. С., Комбинаторная геометрия различных классов вы-

пуклых множеств, Киш., 1978.



ПОКРЫТИЯ И УПАКОВКИ - комбинаторные конфигурации, связанные с многозначным отображением





\end{document}